\chapter{Tuples, Dictionaries \& Sets}

Python provides specialized container structures beyond lists, each tailored for performance, immutability, or fast lookup operations: \textbf{Tuples} (immutable sequences), \textbf{Dictionaries} (key-value hash maps), and \textbf{Sets} (unordered unique collections).

\section{Tuples: Immutable Sequences}

A \textbf{Tuple} is an ordered, immutable sequence of elements defined using parentheses \lstinline|()|.

\begin{definitionbox}{Tuple Properties}
\begin{lstlisting}
point = (10, 20, 30)
print(point[0])    # 10

# Tuple Unpacking
x, y, z = point    # x=10, y=20, z=30

# Immutability: point[0] = 99 raises a TypeError!
\end{lstlisting}
\end{definitionbox}

\textbf{Data Science Relevance:} Tuples guarantee data integrity for fixed coordinates, dataset shapes (\lstinline|df.shape| returns \lstinline|(rows, cols)|), and function return values.

\begin{videocard}{IITM BS Lecture: Tuples Explained}{https://www.youtube.com/watch?v=z-n9yQaWr7o}
Official lecture explaining tuple immutability, packing, unpacking, and hashability.
\end{videocard}

\section{Dictionaries: Key-Value Mapping}

A \textbf{Dictionary} store elements in \lstinline|{key: value}| pairs. Dictionaries provide $\mathcal{O}(1)$ average-time complexity for retrieval by key.

\begin{theorembox}{Dictionary Rules}
\begin{itemize}
    \item \textbf{Keys must be unique and immutable (hashable):} Strings, numbers, and tuples can be keys; lists cannot.
    \item \textbf{Values can be anything:} Lists, nested dictionaries, functions, objects.
\end{itemize}
\end{theorembox}

\begin{examplebox}{Word Frequency Counting}
\begin{lstlisting}
text = "apple banana apple orange banana apple"
words = text.split()

counts = {}
for w in words:
    counts[w] = counts.get(w, 0) + 1

print(counts)  # Output: {'apple': 3, 'banana': 2, 'orange': 1}
\end{lstlisting}
\end{examplebox}

\textbf{Data Science Relevance:} Dictionaries are used to represent JSON API responses, category-to-index mappings in Natural Language Processing (vocabularies), and model hyperparameter configuration dicts.

\begin{videocard}{IITM BS Lecture: Dictionaries for Efficient Storage}{https://www.youtube.com/watch?v=X8Nj5cxaP9E}
Official lecture detailing dictionary creation, \lstinline|.keys()|, \lstinline|.values()|, \lstinline|.items()|, and \lstinline|.get()| methods.
\end{videocard}

\section{Sets: Unordered Unique Collections}

A \textbf{Set} is an unordered collection of unique elements defined using \lstinline|{}| or \lstinline|set()|.

\begin{formulabox}{Set Operations}
\begin{itemize}
    \item \textbf{Union (\lstinline|A | B|):} Elements in $A$, $B$, or both.
    \item \textbf{Intersection (\lstinline|A \& B|):} Elements in both $A$ and $B$.
    \item \textbf{Difference (\lstinline|A - B|):} Elements in $A$ but not in $B$.
    \item \textbf{Symmetric Difference (\lstinline|A \textasciicircum{} B|):} Elements in $A$ or $B$ but not both.
\end{itemize}
\end{formulabox}

\begin{videocard}{IITM BS Lecture: Sets and Set Operations}{https://www.youtube.com/watch?v=qoV4tdDD8zE}
Official lecture explaining set mathematical operations, duplicate removal, and membership testing speed.
\end{videocard}

\newpage
\section{Practice Exercises}
\textit{Detailed step-by-step solutions are available in the Appendix.}

\begin{enumerate}
    \item \textbf{Tuple Operations:} Given \lstinline|student = ("Alice", 21, "Data Science", 3.8)|:
    \begin{enumerate}
        \item Unpack this tuple into variables \lstinline|name|, \lstinline|age|, \lstinline|major|, and \lstinline|gpa|.
        \item What happens if you try to modify \lstinline|student[1] = 22|?
    \end{enumerate}

    \item \textbf{Dictionary Lookup \& Updating:} Create a dictionary \lstinline|stock_prices = {"AAPL": 175.5, "GOOGL": 140.2, "MSFT": 380.0}|:
    \begin{enumerate}
        \item Update the price of \lstinline|"AAPL"| to \lstinline|180.0|.
        \item Add a new key \lstinline|"AMZN"| with price \lstinline|150.0|.
        \item Safely fetch the price of \lstinline|"NVDA"| using \lstinline|.get()| without raising a \lstinline|KeyError|, returning \lstinline|0.0| if missing.
    \end{enumerate}

    \item \textbf{Set Membership \& Deduplication:} Given a raw customer ID list with duplicates \lstinline|ids = [101, 102, 101, 105, 102, 108]|:
    \begin{enumerate}
        \item Convert the list to a set to remove all duplicates.
        \item Check if customer \lstinline|105| is in the set.
    \end{enumerate}

    \item \textbf{Set Algebra:} Given two user cohorts:
    $$A = \{\text{"Alice"}, \text{"Bob"}, \text{"Charlie"}\}, \quad B = \{\text{"Bob"}, \text{"David"}, \text{"Alice"}\}$$
    \begin{enumerate}
        \item Find the intersection of $A$ and $B$.
        \item Find the unique users across both cohorts (union).
    \end{enumerate}

    \item \textbf{Data Science Application (Categorical Encoding Map):} In preprocessing categorical labels into numeric indices for Machine Learning:
    Given \lstinline|labels = ["cat", "dog", "cat", "bird", "dog"]|, write a Python snippet to extract unique labels into a dictionary \lstinline|label2idx| mapping each unique label to a unique integer index (e.g. \lstinline|{"cat": 0, "dog": 1, "bird": 2}|).
\end{enumerate}